\documentclass[11pt,a4paper]{article}
\usepackage[T1]{fontenc}
\usepackage[utf8]{inputenc}
\usepackage{lmodern}
\usepackage[margin=26mm,headheight=14pt]{geometry}
\usepackage{amsmath,amssymb,amsthm,mathtools}
\usepackage{booktabs,longtable,array,microtype}
\usepackage[dvipsnames]{xcolor}
\usepackage{enumitem,xurl,hyperref,listings,fancyhdr}
\hypersetup{colorlinks=true,linkcolor=MidnightBlue,citecolor=MidnightBlue,
 urlcolor=MidnightBlue,pdftitle={Reanchored sign constraints and adaptive Pell carriers},
 pdfauthor={AI-assisted research contribution for ProveIt}}
\urlstyle{same}
\setlist{itemsep=3pt,topsep=5pt}
\setlength{\parskip}{4pt}
\setlength{\parindent}{0pt}
\allowdisplaybreaks[2]
\emergencystretch=2em
\pagestyle{fancy}
\fancyhf{}
\fancyhead[L]{\small Reanchored sign constraints and adaptive Pell carriers}
\fancyhead[R]{\small ProveIt}
\fancyfoot[C]{\thepage}
\renewcommand{\headrulewidth}{0.3pt}
\newtheorem{theorem}{Theorem}[section]
\newtheorem{lemma}[theorem]{Lemma}
\newtheorem{proposition}[theorem]{Proposition}
\newtheorem{corollary}[theorem]{Corollary}
\theoremstyle{definition}
\newtheorem{remark}[theorem]{Remark}
\newcommand{\Z}{\mathbb Z}
\newcommand{\N}{\mathbb N}
\newcommand{\Zp}{\mathbb Z_{>0}}
\newcommand{\Con}{\operatorname{Con}(\mathrm{ZFC})}
\newcommand{\ZFCRef}{\operatorname{Ref}_{\mathrm{ZFC}}}
\newcommand{\Pos}{\operatorname{Pos}}
\newcommand{\MC}{M_C}
\newcommand{\MF}{M_F}
\newcommand{\Dstar}{D_{\!*}}
\newcommand{\Baux}{B_{\mathrm a}}
\newcommand{\Kaux}{K_{\mathrm a}}
\newcommand{\qzero}{q_{\!0}}
\newcommand{\kzero}{k_{\!0}}
\newcommand{\Rzero}{R_{\!0}}
\newcommand{\Bzero}{B_{\!0}}
\newcommand{\indeg}[2]{[#1]_{#2}}
\newcommand{\src}[1]{{\urlstyle{tt}\nolinkurl{#1}}}
\lstset{basicstyle=\ttfamily\footnotesize,breaklines=true,columns=fullflexible,
 keepspaces=true,frame=single,framerule=0.2pt,rulecolor=\color{gray}}
\begin{document}
\hypersetup{pageanchor=false}
\begin{titlepage}
\vspace*{12mm}
{\color{MidnightBlue}\large EXPLICIT DIOPHANTINE CONSISTENCY EQUATIONS\par}
\vspace{9mm}
{\LARGE\bfseries Reanchored sign constraints\\[3pt]
and adaptive Pell carriers\par}
\vspace{8mm}
{\large A sextic in 34 unrestricted integers,\\
a degree-30 equation in 13 positive integers,\\
and a degree-110 equation in 12 positive integers\par}
\vspace{10mm}
{\large Research contribution prepared for the ProveIt project\par}
{23 October 2026\par}
\vspace{8mm}
\begin{abstract}
We give three explicit reductions of the current ProveIt Diophantine
consistency constructions. First, reanchoring common-sign constraints
allows a monic transport elimination while keeping every residual of degree at most three.
Seven coordinates are removed from a 41-integer quartic skeleton,
yielding a degree-six polynomial in 34 unrestricted integers, compared
with the collection's 35-variable sextic. A separately proved complete
coefficient adapter reduces the four numeral payloads from 532,479 to
295,388 raw bytes. Second, replacing the main Pell carrier
\(q^5,q^4X\) by \(q^2Y,qXY\) lowers a thirteen-positive-variable
equation from degree 34 to degree 30, and from 93 to 92 binary
operations. Third, centering both root expressions in an auxiliary
norm lowers a twelve-positive-variable equation from degree 118 to
degree 110, at the cost of one operation.

The new reductions preserve existence of a finite syntactic ZFC
refutation, by composition with explicitly identified source
theorems. Thus ZFC proves their unsolvability equivalent to
\(\Con\). The archive supplies literal polynomials, exact coefficient
payloads, proofs, independent algebra checks, and provenance.
These are local improvements, not a proof of global simplicity or
a new end-to-end proof-assistant certificate. A further appendix
proves reusable Pell and syntax-encoding lemmas; the proposed
direct syntax route remains an incomplete alternative.
\end{abstract}
\vfill
{\small Baseline: ProveIt commit\\
\src{9cca523d9f27ad39d52993f4065cbd7e58f8bdd4}.\\
Unknown counts exclude fixed numerals and polynomial abbreviations.}
\end{titlepage}
\pagenumbering{roman}
\hypersetup{pageanchor=true}
\tableofcontents
\clearpage
\pagenumbering{arabic}

\section{Results and the exact logical claim}
\label{sec:results}

The principal equation is the sum of the squares of the nineteen
residuals of degree at most three in Section~\ref{sec:s34}. It has exactly 34 unrestricted
integer unknowns and exact degree six. The smaller-dimensional
companions use positive integers and have higher degrees. These
are different useful simplicity criteria; none of these points
dominates all the others.

\begin{table}[htbp]
\centering\small
\begin{tabular}{@{}llrrr@{}}
\toprule
Construction & Domain & Unknowns & Degree & Binary operations\\
\midrule
Source 62 sextic & \(\Z\)&35&6&not compared\\
New \(S_{34}\)&\(\Z\)&\textbf{34}&\textbf{6}&not optimized\\
\addlinespace
Source 61 \(P_{13}\)&\(\Zp\)&13&34&93\\
Adaptive \(P_{13}\)&\(\Zp\)&\textbf{13}&\textbf{30}&\textbf{92}\\
\addlinespace
Source 61 product&\(\Zp\)&12&118&104\\
Centered \(P_{12}\)&\(\Zp\)&\textbf{12}&\textbf{110}&105\\
Centered SOS companion&\(\Zp\)&12&164&104\\
\bottomrule
\end{tabular}
\caption{Comparison with the inspected repository. An operation is
a live binary \(+,-,\times\) gate, including multiplication by a
numeral and each squaring. Register reuse and numeral leaves are free.}
\label{tab:results}
\end{table}

The previous twelve-variable SOS companion has degree 172 and
103 operations. Thus the new twelve-variable points exchange
one operation for a lower degree; they do not dominate the old
points in every metric. The thirteen-variable point improves
both degree and operation count.

The baseline already contains nine-positive-variable and
ten-unrestricted-integer equations of much larger degree, and a
quartic in 41 unrestricted integers. In particular, 34 is not a
new fewest-variable bound for arbitrary-degree equations.
The wave-9 nine-positive family includes degree \(1\,494\,006\)
at 487 operations and other tradeoffs \cite{wave67,wave66}.
The repository's closed executable nine-natural-variable Lean
theorem is also a distinct achievement whose proof status must
not be confused with that of these native literal equations
\cite{status}.

\subsection{Consistency and provability}

Fix the repository's recursive predicate \(\ZFCRef(p)\) for a finite
derivation of falsity from its specified ZFC axiom predicate, and
write
\[
 \Con\;:\!\!\iff\;\neg\exists p\in\N\ \ZFCRef(p).
\]
The claimed arithmetical statements, with all coefficient numerals
fixed as specified in Section~\ref{sec:coefficients}, are
\begin{align}
 \mathrm{ZFC}&\vdash
 \bigl(\neg\exists\mathbf x\in\Z^{34}\;S_{34}(\mathbf x)=0\bigr)
       \longleftrightarrow\Con,                 \label{eq:claim34}\\
 \mathrm{ZFC}&\vdash
 \bigl(\neg\exists\mathbf x\in\Zp^{13}\;P_{13}(\mathbf x)=0\bigr)
       \longleftrightarrow\Con,                 \label{eq:claim13}\\
 \mathrm{ZFC}&\vdash
 \bigl(\neg\exists\mathbf x\in\Zp^{12}\;P_{12}(\mathbf x)=0\bigr)
       \longleftrightarrow\Con.                 \label{eq:claim12}
\end{align}
They follow by proving existence of a zero equivalent to existence
of a refutation and negating. No consistency assumption enters a
coefficient definition or witness construction.

Here ZFC-provability is justified by ordinary mathematics
formalizable in ZFC, composed with the cited source arithmetic
and compiler theorems. This contribution proves the new
reductions, not the entire upstream compiler again. It supplies
neither a new end-to-end Lean/Rocq certificate nor a derivation
encoded in the repository's internal ZFC calculus. Finite scripts
audit identities and data; they cannot prove consistency by
testing candidate solutions.

\subsection{The fixed ports}

The four ports \(\beta,K_0,\MC,\MF\) are fixed integer numerals,
not unknowns. The archive supplies losslessly compressed,
minimal unsigned little-endian encodings and deterministic
finite reconstruction, with raw-byte hashes.

The directly inherited \(S_{34}\) instance uses the entire
source-34 tuple with \(\beta=2^{1064921}-1\).
The preferred compact instance, and both positive companions,
use the entire compact source-44 tuple with
\(\beta=2^{590069}-1\). A full adapter for the compact integer
instance is proved in Section~\ref{subsec:compact-adapter}.
In both tuples \(\MF\) is already the shifted high mask.
One must never mix isolated ports from these different tuples.

\section{Optimizing common-sign constraints}
\label{sec:sign}

For integers \(s,t,v\), define
\[
 \Pos(s,t,v)=1+s^2+t^2+v(v+1).                  \label{eq:pos}
\]
This is always positive. Conversely, every positive integer \(m\)
has this form. Legendre's three-square theorem applies to
\(4m-3\), which is \(1\) or \(5\bmod8\). Exactly one square in
such a representation is odd modulo four. Writing it as
\((2v+1)^2\) and the other two as \((2s)^2,(2t)^2\) gives the
formula after division by four.

\begin{lemma}[Changing the sign-constraint tree]
\label{lem:tree}
Let \(x_1,\ldots,x_m\) be integers, \(m\ge2\), and \(T\) any
spanning tree on their indices. The inequalities
\[
 x_ix_j>0\qquad(\{i,j\}\in E(T))
\]
are equivalent to all the \(x_i\) having one common nonzero sign.
Replacing \(T\) by another spanning tree preserves existential
solvability when each edge inequality is represented by its own
fresh ternary \(\Pos\) witness.
\end{lemma}
\begin{proof}
A positive edge product gives its endpoints the same nonzero
sign, and connectivity propagates that sign to every vertex.
The converse is immediate. Surjectivity of \(\Pos\) supplies
the witnesses in both directions.
\end{proof}

Witness triples generally have to be reselected: their represented
products change. This is an existential equivalence, not an
identity with every old coordinate left fixed.

\begin{proposition}[Best tree for fixed polynomial degrees]
\label{prop:tree-degree}
If the nonzero expressions substituted for \(x_i\) have degrees
\(d_i\), then
\[
 \min_T\max_{\{i,j\}\in E(T)}(d_i+d_j)
       =\min_i d_i+\max_i d_i.
\]
A star centered at a minimum-degree expression attains this
bound. The resulting positivity residuals have degree at most
\(\max(2,\min d_i+\max d_i)\).
\end{proposition}
\begin{proof}
Every tree gives a maximum-degree vertex a neighbor of degree
at least \(\min d_i\). The proposed star attains this lower
bound. The positivity polynomial has degree two.
\end{proof}

After the transport substitution below, \(F,\alpha,C\) have
degrees \(2,2,1\). The old edges \(F\alpha,FC\) have degrees
four and three. The new edges \(C\alpha,CF\) both have degree
three. Reanchoring the guards therefore permits a sextic sum
of squares. This elementary principle is a useful optimization
method, with no claim of historical priority.

\section{A sextic in 34 unrestricted integers}
\label{sec:s34}

\subsection{The complete polynomial}

Supply the following 25 unrestricted integers,
\[
\begin{gathered}
 J,C,z,f,h,r,v,w,t,\eta,\zeta,y,\gamma,o,j,\\
 N,E,c,a,D,V,I,\Baux,\Delta,u,
\end{gathered}
\]
and three further triples \(\mathbf p,\mathbf s,\mathbf e\in\Z^3\).
There are \(25+9=34\) unknowns. In the literal files \(v,w,t\)
are named \src{v_scale}, \src{w_scale}, \src{tau_gap}, and the
triples are \src{CA_pos*}, \src{FC_pos*}, \src{eta_zeta_pos*}.

With one complete fixed tuple, put
\begin{align}
 q&=\beta J+1,&k&=\eta+\zeta,&H&=4a+3, \notag\\
 \widehat F&=(K_0+4w)C-z(q-1),&
 \widehat\alpha&=q-\widehat F-C, \notag\\
 \widehat R&=(q-\widehat F)(N-q)-C(q^2-1)
                         +(\MC+q\MF)J.          \label{eq:s34abbr}
\end{align}
These are abbreviations, not extra coordinates. In particular,
\(N\) remains an independent unknown; substituting \(q^3\) into
\(\widehat R\) prematurely would raise its degree.

Define nineteen residuals:
\begin{align}
 L_1&=N-q^3, \notag\\
 L_2&=2u^2-4wq, \notag\\
 L_3&=a-E-(2v+1)N, \notag\\
 L_4&=2u^2(E-a)+E, \notag\\
 L_5&=c-\eta+(E-a)k, \notag\\
 L_6&=t^2+(2t-k)E(c-\eta)-1, \notag\\
 L_7&=D-2u^2-ac-\gamma H, \notag\\
 L_8&=\Delta-a^2-4a-3, \notag\\
 L_9&=D^2-\Delta c^2-1, \notag\\
 L_{10}&=k-1-hE-\widehat R, \notag\\
 L_{11}&=V+c-fo, \notag\\
 L_{12}&=V-cj-hE+k-1, \notag\\
 L_{13}&=I-(2r+1)c^2, \notag\\
 L_{14}&=\Baux-I^2, \notag\\
 L_{15}&=\Baux-\Delta(f^2-1), \notag\\
 L_{16}&=\Baux(V^2-y^2)+y^2-1, \notag\\
 L_{17}&=C\widehat\alpha-\Pos(\mathbf p), \notag\\
 L_{18}&=\widehat FC-\Pos(\mathbf s), \notag\\
 L_{19}&=\eta\zeta-\Pos(\mathbf e).              \label{eq:19rows}
\end{align}
The single Diophantine equation is
\[
 \boxed{\displaystyle S_{34}=\sum_{\nu=1}^{19}L_\nu^2=0,
           \qquad\text{all 34 inputs in }\Z.}   \tag{S34}
\]
There are no implicit divisions, variable exponents, or additional
root witnesses in this definition.

\subsection{Exact local equivalence}

Let \(Q_{41}\) denote the sum of the 26 squared residuals in the
pinned source-64 \src{integer41/quartic41.json}. Its byte-exact
copy is in \src{integer34/upstream/}; the Git blob is
\src{521f16a310f49e736ea8ffac3fe13847ff71adcb}.

\begin{theorem}
\label{thm:s34}
For any fixed integer values of the four coefficient ports, the two
polynomial skeletons satisfy
\[
 \exists\mathbf x\in\Z^{34}\ S_{34}(\mathbf x)=0
 \quad\Longleftrightarrow\quad
 \exists\mathbf y\in\Z^{41}\ Q_{41}(\mathbf y)=0.
\]
\end{theorem}
\begin{proof}
Replace the source coordinate \(\alpha\) by \(C=q-F-\alpha\),
whose inverse is \(\alpha=q-F-C\). Use Lemma~\ref{lem:tree}
to replace \(F\alpha>0,FC>0\) by \(C\alpha>0,FC>0\),
reselecting the first witness triple.

The transport row becomes monic in \(F\):
\[
 (K_0+4w)C-z(q-1)-F=0.
\]
Eliminate \(F=\widehat F\). The two new guards
\(C(q-\widehat F-C)>0,\widehat FC>0\) are cubic; the old
first guard would have become quartic.

Six other monic rows eliminate
\[
 L=q^2-1,\quad X=2u^2,\quad C_2=c^2,\quad F_2=f^2,\quad
 g=y^2-V^2,\quad P=k+E(c-\eta).                 \label{eq:eliminations}
\]
Retain \(N\), so the source \(N=q(L+1)\) becomes \(N=q^3\).
Change the first Pell-root coordinate by the polynomial
automorphism \(\tau=t+E(c-\eta)\). Its source row
\(\tau^2-P(P-k)-1=0\) becomes
\[
 t^2+(2t-k)E(c-\eta)-1=0;
\]
the nominal quartic terms cancel.

All remaining rows are exact substitutions. The packing identity is
\[
 \bigl((q-1)(q-F)+\alpha\bigr)(q^2-1)
   =(q-F)(N-q)-C(q^2-1)
\]
when \(\alpha=q-F-C\) and \(N=q^3\).
This proves that every source zero projects to a new zero.

Conversely, a zero of \(S_{34}\) makes every residual zero.
Restore \(F=\widehat F,\alpha=\widehat\alpha,\tau=t+E(c-\eta)\)
and the six coordinates in \eqref{eq:eliminations}.
The guards give \(C\alpha>0,FC>0\), hence \(F\alpha>0\).
Choose a source positivity triple for the latter product.
Every eliminated row holds by definition, and all retained rows
hold by the stated identities. This gives a \(Q_{41}\) zero.
\end{proof}

The complement and gap changes preserve dimension. Seven rows
and seven coordinates are eliminated, so \(41-7=34\) and
\(26-7=19\). Every remaining variable occurs for each of the two
specified coefficient tuples.

\begin{corollary}
The degree of \(S_{34}\) is exactly six for every coefficient
specialization. With the original source-34 tuple,
\eqref{eq:claim34} follows from the source \(Q_{41}\) theorem.
\end{corollary}
\begin{proof}
Every \(L_\nu\) has degree at most three. Only \(L_{13}\)
involves \(r\), and
\[
 [r^2c^4]S_{34}=4.
\]
This proves exact degree six. The source theorem
\(\exists Q_{41}=0\iff\exists p\,\ZFCRef(p)\)
\cite{q41,q44}, composed with Theorem~\ref{thm:s34}, proves
the consistency equivalence.
\end{proof}

\subsection{A complete compact-tuple adapter}
\label{subsec:compact-adapter}

The same skeleton also accepts the \emph{entire} compact
source-44 tuple, as a separately identified instance. This is
not inferred from small numerical bounds alone. Its semantic
adapter and canonical parity promises are explicitly present
in source 44 and source 36 \cite{compact,parity}.

\begin{theorem}
With the compact tuple defined in Section~\ref{sec:coefficients},
\[
 \exists\mathbf x\in\Z^{34}\ S_{34}(\mathbf x)=0
 \quad\Longleftrightarrow\quad
 \exists p\in\N\ \ZFCRef(p).
\]
\end{theorem}
\begin{proof}
Restore the \(Q_{41}\) skeleton using Theorem~\ref{thm:s34}.
Write \(Q=|q|\), \(X=4qw=2u^2\), \(Y=(2v+1)q^3\),
and let the ordinary strong multiplier be \(i=2r+1\).
The actual compact coefficients satisfy
\[
 \beta=2^d-1,\quad d=590069\ge3,\quad
 \beta\ge17,\quad0<\MC<\beta<\MF<2\beta,\quad K_0\ge0.
\]
The common-sign simplex therefore gives the same sign-free
source-64 bounds
\[
 Q\ge16,\qquad Q(Q^2-1)<R<Q^4-2,\qquad R>3Q+3.
\]
The main norm modulo \(4a+3\) excludes \(X=0\), so
\(X\ge4Q\) and \(|XY|\ge4Q^4>3R+1\).
If both \(\eta,\zeta\) were negative, the first Pell
congruence would imply \(n>R\), while the ordinary strong
and signed auxiliary lemmas imply \(p=R\) with \(p>n\);
this is impossible. Thus \(\eta,\zeta>0\).

The rest of the source-64 signed Pell argument uses only these
bounds and the literal norm and quotient rows. It rules out
\(Y<0\), recovers the exact two ranks, selects the positive
main root, and proves \(X=2^R\). None of these steps depends
on the old block width or old native geometry.

Now \(q\mid2^R\). The elementary Mersenne sign lemma
\(-2^t\not\equiv1\bmod(2^d-1)\), \(d\ge3\), excludes
negative \(q\). To prove that lemma, reduce \(t\) modulo \(d\):
the residue \(2^r\), \(0\le r<d\), lies below \(2^d-2\).
Thus \(q=2^{t_{\rm rad}}>0\), and \(q\equiv1\bmod(2^d-1)\)
also gives \(d\mid t_{\rm rad}\). Hence \(q=B^{N_{\rm tab}}\) for
the compact base \(B=2^d\).

The source half-binomial and inverse-population stages restore
the compact native masks and the shifted transport
\[
 (K_0+w_{\rm source})C=F+z(q-1),\qquad
 w_{\rm source}=X/q.
\]
The rank bound gives \(z>0\); the unshifted convention uses
\(z_{\rm old}=z+1\).
The full source-44 native proof, Sections 5--7, now applies
to its own \(D_C,D_R\), masks, and rotation geometry.
In particular its \(D_R=\rho^{42}\) is used, not the
source-34 \(D_R=\rho\). It decodes a genuine accepting blank
K215 tableau. The inherited literal-machine calibration gives
a finite ZFC refutation.

For completeness, start with such a refutation. Source 44
constructs a padded tableau with
\(N_{\rm tab}=581197^\kappa>1\) odd and
\(t_{\rm rad}=dN_{\rm tab}\) odd.
Its actual-rank controller tunes the rotation, not merely a
surrogate rank. Source 44 Section 6 explicitly preserves
\[
 \operatorname{pc}(R)=3t_{\rm rad}+2,\qquad
 R\equiv3\pmod4,\qquad R>t_{\rm rad},
\]
and the square- and odd-scale promises of source 36 Section 6.
The latter spells out \(v_2(Y)=3t_{\rm rad}\), hence
\(s=Y/q^3\) is positive odd, and the canonical ordinary
strong multiplier \(i\) is positive odd.
Since \(R-t_{\rm rad}\) is positive and even,
\(w_{\rm source}=2^{R-t_{\rm rad}}\) is divisible by four.
Choose
\[
 w=w_{\rm source}/4,\quad
 u=2^{(R-1)/2},\quad v=(s-1)/2,\quad r=(i-1)/2.
\]
The remaining positive canonical coordinates and their exposed
quotients satisfy the \(Q_{41}\) skeleton. Supply the three
ternary positivity triples, then apply Theorem~\ref{thm:s34}.
\end{proof}

The complete interface-by-interface proof and its source blobs
are recorded in \src{integer34/COMPACT_S34_ADAPTER.md}.
The two literal files remain separate:
\src{sextic34.json} retains its source-34 binding, while
\src{sextic34_compact.json} selects the whole compact tuple.
The compact instance uses 295,388 raw coefficient bytes,
against 532,479: a reduction of 237,091 bytes, or \(44.5259\%\).
The degree certificate is independent of the ports and remains
\([r^2c^4]=4\).

\section{An adaptive carrier in thirteen positive unknowns}
\label{sec:p13}

\subsection{Complete polynomial}

Use the compact coefficient tuple. Supply exactly
\[
 j,F,\alpha,f,i,T,X,Y,r,\eta,\zeta,u,\gamma\in\Zp.
\]
Every other symbol is a polynomial abbreviation:
\begin{align}
 J&=3j+1,\qquad q=\beta J+1,\qquad C=q-F-\alpha, \notag\\
 E&=XY,\qquad k=\eta+\zeta,\qquad c=kY+\eta, \notag\\
 a&=E+Y,\qquad H=4a+3,\qquad\Delta=a^2+H, \notag\\
 \ell&=EkY,\qquad b_0=X+ac,\qquad I=ic^2, \notag\\
 R&=\bigl((q-1)(q-F)+\alpha\bigr)(q^2-1)
                    +(\MC+q\MF)J, \notag\\
 r_0&=k-R-1,\qquad U=q^2Y, \notag\\
 \Dstar&=U b_0+H(U\gamma+qXY-r_0), \notag\\
 W&=(q-1)k+1,\qquad M_0=q^3(q-1), \notag\\
 B_t&=(qK_0+X)C-qF, \notag\\
 Z&=M_0(\ell+k+r)+(q-1)Y-q^3WB_t, \notag\\
 \Kaux&=\Delta(f^2-1),\qquad
 V=f^2(cT-R)-c.                                \label{eq:p13registers}
\end{align}
The four residuals are
\begin{align}
 A_0&=Z^2-M_0^2\bigl(\ell(\ell+k)+1\bigr), \notag\\
 A_1&=\Dstar^2-U^2(1+\Delta c^2), \notag\\
 A_{\mathrm a}&=(V+u)^2-\Kaux u(2V+u)-1, \notag\\
 A_{\mathrm s}&=f^2-\Delta I^2-1.                \label{eq:p13residuals}
\end{align}
Set
\[
 \boxed{P_{13}=A_0^2+A_1^2+A_{\mathrm a}^2+A_{\mathrm s}^2=0,
           \qquad\text{all 13 inputs in }\Zp.}   \tag{P13}
\]
The mathematical change from source 61 is precisely
\[
 \boxed{U=q^5,\ q^4X
       \quad\longmapsto\quad U=q^2Y,\ qXY}
\]
inside the main carrier. Shared registers in the literal graph
provide the stated operation count.

\subsection{Inherited Pell and compiler interfaces}

For \(A\ge2\), define the ordinary Pell sequences by
\[
 \chi_A(n)+\psi_A(n)\sqrt{A^2-1}
             =(A+\sqrt{A^2-1})^n .
\]
We use Pell classification, composition, strong divisibility, and
recurrence congruences. The less elementary source interface is
the signed auxiliary-index recovery theorem
\cite{p14,p13old,signed}.

\begin{lemma}[Source signed auxiliary-index interface]
\label{lem:auxiliary}
Let \(A\ge2\), \(\Delta=A^2-1\),
\(c=\psi_A(p)\), \(f=\chi_A(m)\), with \(c\mid m\), \(m>2p\),
and \(0<p,R<c/2\). Suppose \(M>1\), \(c\mid M\), and
\(M^2=\Delta(f^2-1)\). An integer pair \(V,y\) in the nontrivial
branch \(|V|>1\), satisfying
\[
 M^2V^2-(M^2-1)y^2=1,\qquad
 V\equiv-c\pmod f,\qquad V\equiv-R\pmod c,
\]
forces \(p=R\).
\end{lemma}

The cited proof uses odd-index Chebyshev quotient identities and
the plus-congruence step-down theorem. In the application below,
every size, sign, and divisibility hypothesis is checked before
the conclusion is used. No hypothesis \(q\mid X\) belongs to
this interface.

The second inherited interface is semantic. The normalized
positive kernel recovered below, with the whole compact tuple,
has a solution exactly when there is a finite ZFC refutation.
A refutation supplies canonical witnesses with
\begin{gather}
 q=B_{\rm native}^{\,N_{\rm tab}},\quad
 B_{\rm native}=2^{590069},\quad
 J=1+B_{\rm native}+\cdots+B_{\rm native}^{\,N_{\rm tab}-1},
 \notag\\
 N_{\rm tab}=581197^v>1\text{ odd},\quad R\equiv3\pmod4,
 \quad X=qw=2^R,\quad Y=sq^3, \notag\\
 c=\psi_{a+2}(R),\qquad
 k=2\psi_{2XY^2+1}\bigl((R+1)/2\bigr).           \label{eq:canonical}
\end{gather}
Padding permits these choices. The same source supplies the
strong and squared-modulus auxiliary witnesses used for
completeness below \cite{p14,compact}. These are mathematical
interfaces, not empirical conclusions from finite sampling.

\subsection{Soundness in the necessary order}

Suppose all four residuals vanish. We repeatedly use the elementary
fact that a rational number with integer square is an integer:
in lowest terms, its denominator's square divides its numerator's
square, forcing the denominator to be one.

\paragraph{First carrier and recovered signs.}
From \(A_0=0\), rational-root integrality gives \(M_0\mid Z\)
and an integer, possibly signed, \(\tau=Z/M_0\), with
\(\tau^2=\ell(\ell+k)+1\).
Reducing \(M_0\mid Z\) modulo \(q^3\) and modulo \(q-1\)
gives
\[
 q^3\mid Y,\qquad q-1\mid B_t,
\]
since \(\gcd(q,q-1)=1\) and \(W\equiv1\bmod(q-1)\).
Put \(s=Y/q^3>0\), \(t=B_t/(q-1)\). The carrier becomes
\[
 \tau=\ell+k+r+s-Wt.
\]
The norm gives \(|\tau|<\ell+k\), so \(Wt>0\), hence
\(t>0\) and \(B_t>0\). If \(C\le0\), however,
\((qK_0+X)C-qF<0\). Thus \(C>0\), before any use of
\(q\mid X\). The authentic coefficient bounds imply
\[
 q\ge16,\qquad 3q<R<q^4-2.                       \label{eq:Rbounds}
\]
Also \(\beta\equiv J\equiv1\bmod3\), so
\[
 q\equiv2\pmod3,\quad H\equiv3\pmod q,\quad
 \gcd(q,H)=1.                                   \label{eq:coprime}
\]

\paragraph{Positive main root and the correct submodulus.}
We have \(U=q^2Y\ge q^5\), \(r_0<k\), \(c>kY\),
and \(UaY>H\). Therefore
\[
 \Dstar>Uac-Hk>k(UaY-H)>0.
\]
From \(A_1=0\), obtain \(D=\Dstar/U\in\Zp\), with
\[
 D^2-\Delta c^2=1,\qquad
 U\mid H(qXY-r_0).                              \label{eq:adaptive-div}
\]
Now \(q^5\mid U\), while \(q^4\mid qXY\).
Using only the \(q\)-power submodulus and
\eqref{eq:coprime} gives
\[
 q^4\mid r_0.                                   \label{eq:q4r0}
\]
There is no assertion that \(\gcd(U,H)=1\); cancelling
against all of \(U\) at this stage would be unjustified.

\paragraph{Main index before the scale relation.}
Put \(A=a+2\) and \(P_0=2XY^2+1\).
The two Pell classifications give
\[
 |\tau|=\chi_{P_0}(n),\quad k=2\psi_{P_0}(n),\qquad
 D=\chi_A(p),\quad c=\psi_A(p),\qquad n,p\ge1.
\]
Since \(c>1\), \(p\ge2\). If \(n=1\), then \(k=2\);
positivity forces \(\eta=\zeta=1\) and \(c=2Y+1\).
But \(c\ge\psi_A(2)=2A\ge4Y+4\), a contradiction.
Thus \(n\ge2\), and
\[
 c>kY\ge4P_0Y>8XY^3\ge8q^9>2R.
\]
Since \(P_0>A\) and \(c>k\), monotonicity gives \(n<p\).
Pell growth also gives \(c>2p\).

The strong equation gives \(f=\chi_A(m)\), \(I=\psi_A(m)\).
As \(c^2\mid I\), strong divisibility gives \(p\mid m\).
Write \(m=pv\). Composition gives \(I=c\psi_D(v)\);
the congruence \(\psi_D(v)\equiv vD^{v-1}\bmod c\)
and \(\gcd(D,c)=1\) imply \(c\mid v\).
Hence \(c\mid m\), and \(m\ge pc>2p\).

Set \(M=\Delta I\), \(y=|V+u|\). Then \(M^2=\Kaux\),
\(c\mid M\), and the auxiliary residual is
\[
 M^2V^2-(M^2-1)y^2=1.
\]
The exposed congruences are \(V\equiv-c\bmod f\) and
\(V\equiv-R\bmod c\), because \(f^2\equiv1\bmod c\).
Here \(f>2c>4\), excluding \(V=\pm1\); the norm excludes
\(V=0\). Thus every hypothesis of Lemma~\ref{lem:auxiliary}
holds, and \(p=R\). No preceding step assumed \(q\mid X\).

\paragraph{First index and scale recovery.}
The first recurrence gives \(k\equiv2n\bmod4XY^2\).
Its modulus is a multiple of \(q^4\), since \(q^3\mid Y\).
Together with \eqref{eq:q4r0}, this gives
\(q^4\mid2n-R-1\). But \(1\le n<p=R<q^4\), so
\(|2n-R-1|<q^4\). Hence \(2n=R+1\).

Consequently \(r_0=k-2n\) is a positive multiple of \(4XY^2\).
The adaptive modulus now has the decisive property
\[
 \frac{4XY^2}{U}=\frac{4XY}{q^2}\in\Z .
\]
Thus \(U\mid r_0\), and \eqref{eq:adaptive-div}
implies \(U\mid HqXY\), hence \(q\mid HX\).
Equation~\eqref{eq:coprime} gives \(q\mid X\).
The numbers \(w=X/q\) and \(h=r_0/E\) are positive integers.
Also
\[
 \gamma_{\rm old}=\gamma+w-r_0/U
\]
is integral and \(D=X+ac+H\gamma_{\rm old}\).
Since \(D>(a+1)c\) and \(c>X\), it is positive.

\paragraph{Restore the transport and source kernel.}
Now \(B_t=q((K_0+w)C-F)>0\).
Since \(q-1\mid B_t\), a positive integer \(z_{\rm new}\)
satisfies \((K_0+w)C-F=(q-1)z_{\rm new}\).
With \(z_{\rm old}=z_{\rm new}+1\),
\[
 (K_0+w)C+(q-F)-z_{\rm old}(q-1)=1.
\]
Use \(T_{\rm old}=fT>0\) and \(y=|V+u|>0\).
Then \(c(T_{\rm old}f-1)-Rf^2=V\).
All source coordinates and relations are restored.
The inherited kernel soundness gives a finite ZFC refutation.

\subsection{Completeness and a simpler positive gap}

Start with canonical compact refutation witnesses
\eqref{eq:canonical}. Their \(J\) is \(1\bmod3\), giving
the positive integer \(j=(J-1)/3\).
The unchanged first-carrier gap is
\[
 r=Wt+\tau-\ell-k-s,\qquad t=q(z_{\rm old}-1)\ge q.
\]
Source growth gives \(z_{\rm old}\ge2\), \(\tau>\ell\),
and \(k>XY>s\); hence \(r>qW-k-s>0\).
Substitution gives \(Z=M_0\tau\).

For the canonical strong root and \(M=\Delta I\), the
squared-modulus source construction, with \(R\equiv3\bmod4\),
gives
\[
 V_*=\chi_M(R)/M>1,\qquad y_*=\psi_M(R)>V_*,
\]
and \(V_*\equiv-c\bmod f^2,\ V_*\equiv-R\bmod c\).
Therefore
\[
 T=\frac{V_*+c+Rf^2}{cf^2}>0,\qquad u=y_*-V_*>0
\]
are integers and restore the auxiliary equation.

Only the main carrier changes its witness:
\[
 \gamma=\gamma_{\rm old}-w+\frac{r_0}{q^2Y}.
\]
This is integral. Since \(r_0\) is a positive multiple of
\(4XY^2\), and \(X=qw,\ Y=sq^3\),
\[
 \frac{r_0}{q^2Y}\ge\frac{4XY}{q^2}=4wsq^2>w.
\]
Thus \(\gamma>\gamma_{\rm old}>0\).
The predecessor's stronger estimate \(\gamma_{\rm old}>w\)
is unnecessary. Now \(\Dstar=UD\), and all four residuals
vanish. This proves completeness and \eqref{eq:claim13}.

\subsection{Exact degree and arithmetic cost}

Give all supplied variables degree one and ports degree zero.
Put
\[
 \qzero=3\beta j,\quad\kzero=\eta+\zeta,\quad
 \Rzero=\qzero^3(\qzero-F),\quad
 \Bzero=(K_0\qzero+X)(\qzero-F-\alpha)-\qzero F .
\]
Exact expansion gives
\begin{align}
 \indeg{A_0}{15}&=-2\qzero^8\kzero^2XY^2\Bzero,\\
 \indeg{A_1}{13}&=8\qzero^2\kzero X^2Y^4
                         (\qzero^2Y\gamma+\Rzero).
\end{align}
The nominal degree-sixteen terms of \(A_0\) cancel.
The four residual degrees are \(15,13,13,14\).
Only \(A_0^2\) reaches degree 30, giving
\[
 \indeg{P_{13}}{30}
        =4\qzero^{16}\kzero^4X^2Y^4\Bzero^2,\qquad
 [j^{16}\eta^4X^4Y^4\alpha^2]P_{13}=4(3\beta)^{16}>0.
\]
Thus the degree remains exactly 30 with the actual fixed
numerals. The literal graph has 92 live gates:
51 multiplications, 24 additions, 17 subtractions.

\section{A twice-centered norm in twelve positive unknowns}
\label{sec:p12}

\subsection{The complete polynomial}

Supply exactly
\[
 j,F,\alpha,i,T,X,Y,r,\eta,\zeta,n,\gamma\in\Zp.
\]
Use the same outer definitions and \(A_0,A_1\), but
set \(J=6j+1\). The inputs \(f,u\) are absent.
Define
\begin{align}
 I&=ic^2,&\mathcal D&=1+\Delta I^2,&
 M&=\Delta I,&d&=M^2-1, \notag\\
 b&=cT-R,&V&=\mathcal D b-c,&
 \nu&=n+I(b+a+2), \notag\\
 A_*&=M\bigl[d(\nu^2+\mathcal D)+1-V^2\bigr]-2d\nu,
 \notag\\
 B_*&=2d(M\nu-1),&
 \mathcal N&=A_*^2-B_*^2\mathcal D.              \label{eq:centered}
\end{align}
Here \(\mathcal D\) is the polynomial square of the missing
strong root, not the earlier main Pell root.
The single equation is
\[
 \boxed{P_{12}=(1+A_0)(1+A_1)(1+\mathcal N)-1=0,
           \qquad\text{all 12 inputs in }\Zp.}   \tag{P12}
\]
The SOS companion
\[
 P_{12}^{\rm SOS}=A_0^2+A_1^2+\mathcal N^2=0     \tag{P12-SOS}
\]
has the same positive zero set, with the same \(J=6j+1\)
convention.

\subsection{The product is an exact conjunction}

Both \(J,\beta\) are odd, so \(q\) is even; therefore
\(M_0,U\) are even. Hence \(A_0,A_1\) are squares modulo four.
So is \(\mathcal N\), since \(B_*\) is even.
If the three integer factors in (P12) multiply to one, each
is \(\pm1\). A factor \(-1\) would require its residual to
equal \(-2\equiv2\bmod4\), impossible for a square modulo four.
Thus \(P_{12}=0\) exactly when \(A_0=A_1=\mathcal N=0\).

\subsection{Soundness of the centered norm}

Positive inputs give \(c\ge3,\ \mathcal D>c+1,\ M>1,\ d>0\).
The abbreviation \(\nu\) need not be positive.
Nevertheless \(M\nu-1\ne0\) for integer \(\nu\), so \(B_*\ne0\).
At \(\mathcal N=0\), let \(f=A_*/B_*\in\mathbb Q\).
Then \(f^2=\mathcal D\in\Z\), hence \(f\in\Z\).
Put \(t=\nu-f,\ y=Mt-1\). Modulo \(f^2=\mathcal D\),
the exact identities are
\begin{align}
 A_*-B_*f&=Md\,t^2-2dt+M(1-V^2),               \label{eq:centerid1}\\
 M(A_*-B_*f)&=d(Mt-1)^2-M^2V^2+1.             \label{eq:centerid2}
\end{align}
Consequently
\[
 M^2V^2-(M^2-1)y^2=1.                          \label{eq:normlift}
\]
Since \(b\in\Z\) and \(V=\mathcal D b-c\), either
\(b\le0,\ V\le-c<-1\), or \(b\ge1,\ V\ge\mathcal D-c>1\).
Thus \(|V|>1\), and \eqref{eq:normlift} gives
\[
 d(y^2-V^2)=V^2-1>0.
\]
Therefore \(|y|>|V|\) and \(u_{13}=|y|-V>0\).
Take \(j_{13}=2j,\ f_{13}=|f|\), and this \(u_{13}\), retaining
all other inputs. The outer \(J,q,A_0,A_1\) agree; the strong
and auxiliary equations follow from \(f^2=\mathcal D\)
and \eqref{eq:normlift}. Every occurrence of \(f\) in (P13)
is through \(f^2\). We have produced a positive \(P_{13}\) zero.

\subsection{Completeness and the positive centered coordinate}

Canonical compact \(J\) is a repunit of odd length
\(N_{\rm tab}>1\) in the even base \(2^{590069}\equiv2\bmod3\).
Thus \(J\) is odd and \(J\equiv1\bmod3\), hence \(J\equiv1\bmod6\).
The input \(j=(J-1)/6\) is therefore a positive integer.
This proves that the product finalizer's parity restriction
loses no canonical refutation witnesses.

Choose the canonical auxiliary pair
\[
 V=\chi_M(R)/M>0,\qquad y=\psi_M(R)>V.
\]
Since \(R\equiv3\bmod4\), the odd-index recurrence gives
\(y\equiv-1\bmod M^2\).
For the canonical strong root \(f>0\), set
\[
 t=(y+1)/M,\qquad \nu=f+t,\qquad
 n=\nu-I(b+a+2).
\]
These are integral. To prove \(n>0\), put
\(\lambda=(a+2)I\), \(u=y-V>0\).
Using \(\mathcal D-1=MI\) and \(V=\mathcal D b-c\),
\[
 n=f-\lambda+\frac{u+b-c+1}{M}.                 \label{eq:npositive}
\]
We have \(M>\lambda,\ M>c,\ \mathcal D<M^2,\ M\ge3\).
Also \(R>3q\) gives \(R\ge7\).
Pell growth \(\chi_M(R)\ge M^R\) implies
\(V\ge M^{R-1}\ge M^6\), whereas
\[
 \mathcal D(M\lambda+c)<M^2(M^2+M)<2M^4<M^6.
\]
Thus \(b=(V+c)/\mathcal D>M\lambda+c\).
Equation~\eqref{eq:npositive} yields \(n>f>0\).
Substitution in \eqref{eq:centerid1}--\eqref{eq:centerid2}
gives \(A_*=B_*f\), hence \(\mathcal N=0\).
The outer residuals already vanish, proving completeness
and \eqref{eq:claim12}.

This chooses fresh canonical auxiliary witnesses. It does not
claim that every old witness has \(y\equiv-1\bmod M^2\) or a
positive shifted coordinate.

\subsection{Exact cancellation and metrics}

First treat \(a,c,i,b,n\) as independent variables with weights
\((2,2,1,4,1)\), with \(\Delta=a^2+4a+3\).
Complete sparse expansion gives 875 monomials in \(\mathcal N\),
weighted degree 82, and highest form
\[
 \indeg{\mathcal N}{82}=4a^{12}c^{20}i^{10}b^2.
\]
The nominal degree-86 terms in \(A_*^2\) and
\(B_*^2\mathcal D\) cancel exactly.
After composition, put \(\qzero=6\beta j\) and retain
\(\kzero,\Rzero,\Bzero\) as before. Then
\[
 \indeg{\mathcal N}{82}
    =4i^{10}\kzero^{20}X^{12}Y^{32}\Rzero^2
\]
and
\[
 \indeg{P_{12}}{110}
   =-64\qzero^{10}i^{10}\kzero^{23}X^{15}Y^{38}\Bzero
              (\qzero^2Y\gamma+\Rzero)\Rzero^2.
\]
The coefficient
\[
 [j^{18}i^{10}\eta^{23}X^{16}Y^{39}\gamma\alpha F^2]P_{12}
       =64(6\beta)^{18}>0
\]
certifies exact degree 110. The graph has 105 live gates:
58 multiplications, 29 additions, 18 subtractions.
The SOS companion has exact degree 164 and 104 gates.

The symmetric auxiliary divisor \(M=\Delta ic^2\), combined
with the shift \(I(b+a+2)\), cancels the polynomial parts of
both root expressions. Changing the divisor split alone does
not produce the stated improvement.

\section{Exact coefficient numerals and provenance}
\label{sec:coefficients}

\subsection{Complete compact recipe}

The preferred instance uses
\[
 \rho=2^{67},\quad L=8807,\quad d=67L=590069,\quad B=2^d.
\]
Take the ordered 943 tiles
\((w_i,n_i,e_i,s_i)\), \(0\le i<943\), from the included
literal \src{compact_coefficients/native/wang943_tiles.json}.
The order matters: tile zero is Start and tile one is Forbidden.
The canonical serialized tile-list SHA256 is
\begin{lstlisting}
66b55f1d11fca9fefc93b12fcfc8b4cabb75b08e9fd03311f9e894a1c5e5e81f
\end{lstlisting}
The builder reconstructs this finite list from all 430 literal
machine actions, without executing the machine; the independent
auditor begins from the list instead.

Define a 41-bit payload for each tile,
\[
 v_i=w_i+2^8n_i+2^{17}e_i+2^{25}s_i+127\cdot2^{34},
 \qquad c_i=1+2^{10}(2^{41}-v_i).
\]
Let \(\mathcal E\) contain the selector positions \(0,\ldots,942\),
the copy and anchor positions in the next table, and the two
dummy positions \(1063,1064\).

\begin{table}[htbp]
\centering\small
\begin{tabular}{@{}lll@{}}
\toprule
Native position & Payload bit \(j\) & Range\\
\midrule
\(1023+h\)&\(h\)&\(0\le h<8\)\\
\(1031\)&\(34\)&west occupancy\\
\(1013+h\)&\(8+h\)&\(0\le h<9\)\\
\(1022\)&\(35\)&north occupancy\\
\(1003+h\)&\(17+h\)&\(0\le h<8\)\\
\(1011\)&\(36\)&east occupancy\\
\(983+h\)&\(25+h\)&\(0\le h<9\)\\
\(992\)&\(37\)&south occupancy\\
\(1095,2195,3296\)&\(38,39,40\)&three anchors\\
\bottomrule
\end{tabular}
\caption{The compact tuple's finite copy-position map.}
\end{table}

At a copy or anchor position carrying bit \(j\), set
\(c_e=2^{10+j}\); set \(c_{1063}=c_{1064}=0\).
There are 986 positions in \(\mathcal E\).
Put \(\mu=2^{51}-2\) and define the following fixed integers:
\begin{align}
 D_C={}&\rho^{30}+\rho^{62}+\rho^{1100}+\rho^{2201}
       +\sum_{e\in\mathcal E}c_e
                         (\rho^{4393-e}+\rho^{5495-e}),\notag\\
 D_R={}&\rho^{42},\notag\\
 Z_{\rm pad}={}&(2^{65}-1)\sum_{e=8792}^{8804}\rho^e
                         +(2^{19}-1)\rho^{8805},\notag\\
 \MC={}&B-1-\sum_{\substack{e\in\mathcal E\\e\ne1}}\rho^e
                                      -2\rho^{1063},\notag\\
 M_{F,\rm native}={}&4+\mu(\rho^{4393}+\rho^{5495})
        +\sum_{h=0}^{9}\rho^{1013+h}
        +\sum_{h=0}^{8}\rho^{1065+h}\notag\\
       &\hspace{18mm}+\rho^{2195}+\rho^{3296}+Z_{\rm pad}.
                                                        \label{eq:compactrecipe}
\end{align}
The polynomial's four fixed ports are exactly
\[
 \boxed{\beta=B-1,\qquad K_0=D_C+BD_R,\qquad
        \MC=\MC,\qquad \MF=M_{F,\rm native}+B-1.}
                                                        \label{eq:ports}
\]
Thus \eqref{eq:compactrecipe}, the finite tile list, and
\eqref{eq:ports} specify actual integers, not unspecified
program indices. All sums and exponents in this coefficient
recipe are fixed and finite.

\begin{table}[htbp]
\centering\small
\begin{tabular}{@{}llrrr@{}}
\toprule
Polynomial leaf & Payload name & Bits & Raw bytes & Gzip bytes\\
\midrule
\(\beta\)&\src{beta.bin.gz}&590069&73759&108\\
\(K_0\)&\src{K.bin.gz}&592884&74111&8940\\
\(\MC\)&\src{MC.bin.gz}&590069&73759&193\\
\(\MF\)&\src{MF_source.bin.gz}&590070&73759&204\\
\midrule
Total&&&295388&9445\\
\bottomrule
\end{tabular}
\caption{The preferred complete compact tuple. Compression is a
delivery format, not a change in mathematical coefficient height.}
\end{table}

Each payload below \src{compact_coefficients/native/coefficients/}
is gzip of the minimal unsigned little-endian integer bytes.
It is read literally by
\begin{lstlisting}
value = int.from_bytes(gzip.decompress(path.read_bytes()), "little")
\end{lstlisting}
The raw hashes, in \(\beta,K_0,\MC,\MF\) order, are
\begin{lstlisting}
0664a868cbfef9869ea05c40a1c445563f685ff857092594c1c1223bc4b21f9d
c6d523bd32af8dd0345db843cdadf6932bfcad186912a6d112314b60657c554e
6007cabd5f5228618e8fb613f386d8fc3692a24520d6e4cd72af0262aae1ade4
48333b0d30a4edeb8579f44286cb59806c5622cb42b367f5b5f9ee827711c737
\end{lstlisting}
The full upstream binding is
\src{compact_coefficients/native/COEFFICIENT_BINDING.json},
and \src{POSITIVE_PORTS.json} records the polynomial aliases.
Its instance identifier is
\begin{lstlisting}
wang943-complement67-shared-anchor-length8807-prime581197
\end{lstlisting}
The selector, mask, carry, geometry, and actual-rank controller
proofs are in the preserved source-44 \src{native/PROOF.md}.
The controller prime is \(581197\), with step \(581196\).
No coefficient is computed by waiting for a halting computation.

\subsection{The original comparison instance}

For the directly inherited \(S_{34}\) binding,
\(\rho=2^{121}\), \(L=8801\), \(d=1064921\), \(B=2^d\).
Its full finite recipe is preserved in
\src{integer_coefficients/README.md}, alongside the exact
source constructor, independent decoder, and four raw and
compressed payloads. Their total raw size is 532,479 bytes;
the largest port has 1,065,043 bits. The complete tuple digest is
\begin{lstlisting}
b61ee51bdd2d9e4274db603b64cb1d092a8f5feeb83f510433876f21fea2c85f
\end{lstlisting}
The manifest specifies its length-delimited digest convention.
This instance remains useful for the direct
\(Q_{41}\)-to-\(S_{34}\) comparison. The compact instance is
preferred when the four coefficient sizes are part of the
simplicity objective.

\subsection{Pinned evidence}

All retrievals are pinned to commit
\src{9cca523d9f27ad39d52993f4065cbd7e58f8bdd4}.
Historical source pins inside copied files are retained
unchanged; provenance ledgers separately record current
retrieval paths, SHA256, Git blob SHA1, and byte counts.
The compact bundle authenticates 44 copied dependencies;
the original tuple bundle authenticates 35.
The preserved source proofs include native soundness,
canonical completeness, signed Pell recovery, and literal
machine/refutation calibration. Some inherited relative links
assume the original repository layout; the ledgers identify
their exact upstream targets.

\section{Verification and repository integration}
\label{sec:verification}

The archive separates finite algebra checks from unbounded
mathematical proofs. The new reductions above establish
existence equivalence using explicit source theorems.
The independent standard-library scripts establish:

\begin{enumerate}
\item All 26 source-row restoration identities for \(S_{34}\),
including the seven eliminated rows; inverse complement and
gap identities; 34 live inputs; nineteen residuals of degree
at most three; and the exact coefficient \([r^2c^4]=4\).
\item A complete expansion of \(S_{34}\) into 437 monomials
when the four coefficient ports remain independent symbols.
This is not a monomial-minimality assertion after specialization.
\item Exact identities for the positive carriers, graph liveness,
gate ledgers, highest-degree forms, and both finalizers.
The checker reads the literal graphs and imports no builder.
\item Complete expansion of the 875-term centered norm in an
independent sparse polynomial implementation, including the
degree-86 cancellation and the two root-recovery identities.
\item Reconstruction of all actual coefficient numerals from
the pinned finite data, with raw-byte hashes and the original
gzip hashes. No small surrogate coefficients are used.
\item Literal equality of the original and compact sextic
polynomial skeletons, with separately authenticated complete
port bindings and adapter proof identities.
\end{enumerate}

The compact native auditor checks all 430 machine actions,
66,062 basis identities, 38,663 single-bit mutations,
and 87,380 finite complement models. The original tuple
audit checks 69,230 basis identities and all 8,801 digits of
its transport constant. These finite diagnostics supplement,
and do not replace, the universal native proofs.

From the extracted package root, run
\begin{lstlisting}
python3 verify_all.py
python3 -O verify_all.py
latexmk -pdf -interaction=nonstopmode -halt-on-error report.tex
\end{lstlisting}
Python 3.10 or later and ordinary TeX Live packages suffice.
The finite checkers run normally and with optimization;
essential validation does not depend on removable
\src{assert} statements. The top-level receipt records
individual process outcomes and logs. Exact sparse
expansions and checked result JSON files are included.

The new local reductions are suitable for additive integration
as a new proposal directory; the archive changes no existing
repository file. Its \src{README.md} identifies the preferred
equation, directory layout, exact commands, comparison basis,
and proof dependencies. Licenses and upstream notices are
preserved. A manifest hashes the delivered files.

\subsection{A concrete formalization target}

Theorem~\ref{thm:s34} can be formalized independently of the
enormous numerals and computation encoding: prove the
sign-tree lemma, the positivity representation, and the
polynomial substitutions over symbolic integer ports.
That local theorem can later be composed with a formal
source equation. The carrier proof similarly identifies
precise divisibility and positivity lemmas.

Neither this package nor its finite verification closes the
native compiler semantic bridge in a proof assistant, recertifies
the machine's full formal calibration, or supplies an internally
encoded ZFC derivation. The repository's existing distinctions
between those targets remain in force \cite{status}.

\section{Search decisions and remaining work}
\label{sec:limits}

The request called for a simpler equation and a serious attempt
at a new approach. The constructions here improve the inspected
frontier in several concrete metrics. They do not establish
global simplicity, literature-wide priority, or a new
fewest-variable record.

Several tempting shortcuts fail at specific obligations:
\begin{itemize}
\item Removing the \(X=2u^2\) guard leaves a negative-power
branch not justified by the existing computation encoding.
The new sextic keeps that guard.
\item The adaptive carrier cannot cancel \(H\) against all
of \(U=q^2Y\). Its proof first uses the \(q^4\) submodulus,
then the exact Pell index, and only then recovers \(q\mid X\).
\item Weakening \(q^3\mid Y\) to \(q^2\mid Y\) is not
automatically compatible with the native population threshold.
A Pell-jet alternative supplies a useful local direction,
but gives a larger degree and does not restore that interface.
\item Replacing the transport multiplier \((q-1)k+1\) by \(k\)
only proves \((q-1)\mid kB_t\); the necessary cancellation
is unavailable.
\item A direct syntax compiler can omit some well-formedness
tests by interpreting unused codes conservatively, as proved
in the appendix. It still needs an explicit axiom theory,
a bidirectional consistency translation, and global
unbounded-proof packing.
\end{itemize}

The most promising further algebraic task is to optimize the
sign-constraint tree and retained registers jointly. The
\(N\)-versus-\(L\) choice and the guard reanchoring interact:
optimizing either alone misses the sextic reduction.
A more ambitious task is a native compiler whose population
threshold can use the weaker \(q^2\mid Y\) scale without first
restoring the old kernel. Neither task is declared solved
here.

\appendix
% Standalone section for \input into a report using amsmath, amssymb, amsthm.
% This is a research appendix, not a separately compiled report.
\newtheorem{auxsyntaxthm}{Theorem}[section]
\newtheorem{auxsyntaxlem}[auxsyntaxthm]{Lemma}
\newtheorem{auxsyntaxprop}[auxsyntaxthm]{Proposition}

\section{Two reusable arithmetic lemmas}
\label{sec:reusable-arithmetic-lemmas}

This appendix proves two independent tools.  The first permits a direct
equational proof compiler to use a simple injective polynomial numbering
without checking that every term code is in the numbering's image.  The
second recovers an exact power of two from a pair of Pell ranks under
weaker hypotheses than a radix-dependent formulation would suggest.
Neither result, by itself, constructs a new equation for
$\operatorname{Con}(\mathrm{ZFC})$.  No claim of historical priority is made.

\subsection{Injective syntax and conservative extra constants}

Throughout this subsection, $\mathbb N=\{0,1,2,\ldots\}$.  Fix the
first-order equational signature $\Sigma$ consisting of constants
$\epsilon,\iota$, binary function symbols $\mathsf U,\mathsf C$, and unary
function symbols $\mathsf N,\mathsf V$.  A fixed theory $\mathcal E$ is a
set of equations in finitely many variables; its axioms are understood
through substitution instances.  The argument applies to any fixed such
set, and in particular to a finite set.

We use the ordinary equational calculus: substitution instances of the
axioms, reflexivity, symmetry, transitivity, and congruence for each
function symbol.  A proof of a closed equation can be taken to consist
entirely of ground equations.  Indeed, replace every variable occurring
in a given finite proof by the ground constant $\epsilon$.  This preserves
the closed conclusion, sends axiom instances to axiom instances, and
preserves every inference.  This observation concerns a specified
syntactic calculus and requires no semantic completeness theorem.

Define integer polynomials
\begin{align}
 J(a,b)&=(a+b)^2+a+1,\\
 U(a,b)&=4J(a,b),& C(a,b)&=4J(a,b)+1,\\
 N(a)&=4a+2,& V(a)&=4a+3.
\end{align}
The code of $\epsilon$ is $0$, and the code of $\iota$ is $1$.
The code $\langle t\rangle$ of a ground $\Sigma$-term is obtained by
applying the corresponding polynomial at every constructor.

\begin{auxsyntaxlem}[Disjoint constructors and decreasing decoding]
The pairing $J:\mathbb N^2\to\mathbb N$ is injective and satisfies
$J(a,b)>\max(a,b)$.  The four constructor images are pairwise disjoint
and disjoint from $\{0,1\}$.  Every argument of a constructor has a
strictly smaller code than its value.
\end{auxsyntaxlem}
\begin{proof}
Put $s=a+b$.  For a fixed $s$, the value of $J$ lies in
$[s^2+1,s^2+s+1]$.  These intervals are disjoint because
$(s+1)^2+1>s^2+s+1$.  The interval identifies $s$, after which
$a=J-s^2-1$ and $b=s-a$ identify the arguments.  Also
$J(a,b)\ge s^2+1>s$ for $s\ge1$, and the case $s=0$ is immediate.
The four constructors have distinct residues modulo $4$; the two binary
constructors have values at least $4$ and $5$, respectively.  The
remaining assertions follow directly from their formulas.
\end{proof}

Let $\mathcal J$ be the set of natural numbers outside the four
constructor images and $\{0,1\}$.  Enlarge $\Sigma$ by a distinct constant
$\kappa_m$ for each $m\in\mathcal J$, obtaining $\Sigma^+$.  Strong
recursion defines a decoding map
$\delta:\mathbb N\to\operatorname{Term}(\Sigma^+)$: decode a constructor
value using its unique arguments, decode $0,1$ as $\epsilon,\iota$, and
decode $m\in\mathcal J$ as $\kappa_m$.  The preceding lemma proves that
this definition terminates.  It also proves, for example,
\[
 \delta(U(a,b))=\mathsf U(\delta(a),\delta(b)),
 \qquad \delta(N(a))=\mathsf N(\delta(a)),
\]
with analogous identities for the other constructors.  Induction gives
$\delta(\langle t\rangle)=t$ for every original ground term $t$.

For an axiom term $t(x_1,\ldots,x_h)$, let
$\widehat t(n_1,\ldots,n_h)$ be its integer-polynomial evaluation under
these arithmetic constructors.  Arbitrary natural values $n_i$ are
allowed.  Constructor compatibility gives
\[
 \delta\bigl(\widehat t(\boldsymbol n)\bigr)
   =t\bigl(\delta(n_1),\ldots,\delta(n_h)\bigr).
\]
Thus no test that $n_i$ belongs to the image of the original term
numbering is needed.

\begin{auxsyntaxthm}[Conservative interpretation of unused codes]
\label{thm:syntax-retraction}
Let $s=t$ be a closed $\Sigma$-equation.  There exists a finite arithmetic
equational proof of
$(\langle s\rangle,\langle t\rangle)$, using unrestricted natural term
codes, polynomial evaluations of $\mathcal E$-instances, and the encoded
equational rules, if and only if $\mathcal E\vdash s=t$ in the specified
ordinary equational calculus.
\end{auxsyntaxthm}
\begin{proof}
An ordinary ground proof can be encoded line by line, proving the
reverse implication.  Conversely, apply $\delta$ to the two term codes
on every line of an arithmetic proof.  The constructor identities show
that it becomes an ordinary ground proof in the signature $\Sigma^+$.
Every arithmetic axiom instance becomes an actual $\mathcal E$-instance
with ground $\Sigma^+$-terms as substitutions.

Now apply the term homomorphism
$\rho:\operatorname{Term}(\Sigma^+)\to\operatorname{Term}(\Sigma)$
which fixes every original symbol and sends every $\kappa_m$ to
$\epsilon$.  It sends an axiom instance to an axiom instance and
preserves each equational rule.  It fixes the closed terms $s,t$.
Consequently the transformed finite proof is an original-signature
proof of $s=t$.  This is a syntactic retraction argument; it assumes
neither consistency nor arithmetic soundness of any theory.
\end{proof}

\subsubsection*{One explicit polynomial for local inference}

Represent a conclusion by its pair of term codes $(a,b)$ and two premises
by $(c,d)$ and $(e,f)$.  For one additional natural number $z$, put
\begin{align}
 R_0&=(a-b)^2,\\
 T_0&=(c-e)^2+(a-d)^2+(b-f)^2,\\
 B_0&=(a-U(c,e)-z)^2+(b-U(d,f)-z)^2,\\
 W_0&=(a-4c-2-z)^2+(b-4d-2-z)^2,\\
 I(a,b,c,d,e,f,z)&=z^2(z-1)^2+R_0T_0B_0W_0.
 \label{eq:local-inference-poly}
\end{align}
Here $U$ in the two displayed occurrences is the already defined
integer polynomial; these are abbreviations, not extra unknowns.

\begin{auxsyntaxprop}[Exact local rule check]
For natural term codes $a,b,c,d,e,f$, the equation $I=0$ holds precisely
when $z\in\{0,1\}$ and at least one of the following holds:
\begin{enumerate}
 \item $a=b$, the reflexivity rule;
 \item $c=e$, $a=d$, $b=f$, the rule taking $c=d$ and $c=f$ to $d=f$;
 \item $(a,b)=(U(c,e),U(d,f))$ if $z=0$, or
       $(a,b)=(C(c,e),C(d,f))$ if $z=1$;
 \item $(a,b)=(N(c),N(d))$ if $z=0$, or
       $(a,b)=(V(c),V(d))$ if $z=1$.
\end{enumerate}
The total degree of $I$ is exactly $10$.  The variable $z$ is its only
local auxiliary coordinate beyond the six equation ports.
\end{auxsyntaxprop}
\begin{proof}
Every displayed factor and summand is nonnegative.  Thus $I=0$ forces
both $z^2(z-1)^2=0$ and $R_0T_0B_0W_0=0$.  The former is exactly the
conjunction $z\in\{0,1\}$; the latter is exactly the disjunction of
the four rule checks.  Each sum of squares vanishes precisely when all
its summands vanish.  This also proves the converse.

The four factors have degrees $2,2,4,2$, so $\deg I\le10$.
Specializing gives the polynomial identity
\[
 I(t,0,t,0,0,0,0)
 =2t^4\bigl((4t^2+3t+4)^2+16\bigr)
       \bigl((3t+2)^2+4\bigr).
\]
Its leading coefficient is $288$, proving $\deg I=10$.
\end{proof}

The second rule and reflexivity derive symmetry: use the premises
$c=d$ and $c=c$.  They then derive ordinary transitivity by first
reversing $a=b$ and using $b=a$, $b=c$.  Binary congruence of both
arguments includes one-argument congruence by using a reflexive second
premise.  Hence these local rules generate the ordinary ground
equational calculus when supplemented with the axiom instances.

To make the finite-proof convention exact, seed the list by the valid
reflexive equation $(0,0)$.  Each later line is either an axiom instance
or has two earlier premise indices and a natural $z$ satisfying
\eqref{eq:local-inference-poly}.  A one-premise rule may use the seed
as its unused second premise.  This convention excludes circular
justifications, including a purported first inference from itself.
The premise-index restrictions and the unbounded list are part of the
finite-certificate definition; they have \emph{not} been eliminated
into the local polynomial $I$.

\paragraph{What remains to obtain a consistency equation.}
One must supply an explicit finite equational theory, prove the required
bidirectional syntactic consistency translation, materialize its actual
axiom terms, and compile the entire unbounded finite-proof certificate
into a fixed Diophantine equation.  No variable count, degree bound, or
arithmetic-operation improvement for that unconstructed endpoint follows
from the local degree-$10$ calculation.

\subsection{Exact Pell power recovery without an early radix bound}

For an integer $b\ge2$, define the Pell sequences by
\[
 \chi_b(0)=1,\quad \psi_b(0)=0,\qquad
 \begin{pmatrix}\chi_b(m+1)\\\psi_b(m+1)\end{pmatrix}
 =
 \begin{pmatrix}b&b^2-1\\1&b\end{pmatrix}
 \begin{pmatrix}\chi_b(m)\\\psi_b(m)\end{pmatrix}.
\]
In particular, $\psi_b(1)=1$ and
$\psi_b(m+2)=2b\psi_b(m+1)-\psi_b(m)$.

\begin{auxsyntaxlem}[Recurrence bounds]
For every $b\ge2$ and integer $m\ge1$,
\[
 (2b-1)^{m-1}\le\psi_b(m)\le(2b)^{m-1}.
\]
\end{auxsyntaxlem}
\begin{proof}
The sequence $\psi_b(m)$ is nonnegative and strictly increasing from
$m=0$ onward.  This follows inductively from
$\psi_b(0)=0$, $\psi_b(1)=1$ and the second-order recurrence.  For
$m\ge1$, the same recurrence therefore gives
\[
 (2b-1)\psi_b(m)
 \le\psi_b(m+1)
 \le2b\psi_b(m).
\]
Iterating from $\psi_b(1)=1$ proves both bounds.
\end{proof}

\begin{auxsyntaxlem}[Exact quotient identity]
\label{lem:pell-quotient-exact}
Let $a\ge1$, $A=a+2$, and $H=4a+3$.  For every $m\ge0$,
\[
 \chi_A(m)-a\psi_A(m)
 =2^m+H\sum_{j=0}^{m-1}2^{m-1-j}\psi_A(j),
\]
where the empty sum is zero.  In particular,
$\chi_A(m)-a\psi_A(m)\equiv2^m\pmod H$.
\end{auxsyntaxlem}
\begin{proof}
Put $W_m=\chi_A(m)-a\psi_A(m)$ and note that
$A^2-1=a^2+H$.  The defining recurrences give
\[
 W_{m+1}
 =(A-a)\chi_A(m)+(A^2-1-aA)\psi_A(m)
 =2W_m+H\psi_A(m).
\]
Since $W_0=1$, iterating this identity yields the asserted formula.
The sign is that of the positive sequence $\chi_A$; it is not a
statement obtained by replacing a possibly signed Pell coordinate by
its absolute value without further justification.
\end{proof}

\begin{auxsyntaxthm}[Power recovery]
\label{thm:pell-power-recovery-clean}
Let $X,Y$ be positive integers and $n\ge2$.  Put
\[
 R=2n-1,\quad a=(X+1)Y,\quad A=a+2,\quad
 P=2XY^2+1,\quad H=4a+3.
\]
Suppose
\[
 c=\psi_A(R),\qquad k=2\psi_P(n),\qquad
 kY<c<k(Y+1),
\]
and suppose the \emph{positive} main root
$D=\chi_A(R)$ satisfies
\[
 D-ac\equiv X\pmod H.
\]
Then $X=2^R$.  No hypothesis involving a radix $q$, divisibility of $X$
or $Y$ by $q$, or inequalities $X\ge q$, $Y\ge q^3$ is required.
\end{auxsyntaxthm}
\begin{proof}
Both $A$ and $P$ are at least $2$.  Direct calculation gives
\[
 (2A-1)^2-8P
 =4(X-1)^2Y^2+12(X+1)Y+1>0.
\]
Since $R-1=2(n-1)$, the recurrence bounds imply
\begin{align*}
 \frac ck
 &\ge \frac{(2A-1)^{2(n-1)}}{2(2P)^{n-1}}\\
 &=\frac12\left(\frac{(2A-1)^2}{2P}\right)^{n-1}
 >\frac{4^{n-1}}2=2^{R-2}.
\end{align*}
The strict ratio bound $c/k<Y+1$ therefore gives
$Y+1>2^{R-2}$.  Since $R\ge3$ and $Y$ is an integer,
$Y\ge2^{R-2}$.  Consequently
\[
 H=4(X+1)Y+3\ge8Y+3\ge2^{R+1}+3>2^R.
\]
Also $0<X<H$ follows directly from $X,Y\ge1$.

Lemma~\ref{lem:pell-quotient-exact} gives
$D-ac\equiv2^R\pmod H$.  The hypothesis gives the same residue as $X$.
Both $X$ and $2^R$ lie strictly between $0$ and $H$, so they are equal.
\end{proof}

The lower ratio bound $kY<c$ is natural in the intended applications,
although this particular proof uses only the upper bound.  If the source
instead supplies $D=X+ac+H\gamma$, then after power recovery the exact
quotient formula identifies
\[
 \gamma=\sum_{j=0}^{R-1}2^{R-1-j}\psi_A(j)>0,
\]
because $R\ge3$.  Thus the same elementary identity also verifies the
sign of the recovered quotient.  The hypotheses giving the ranks
$R=2n-1$, the positive main root, and the strict ratio remain essential
inputs to this lemma; they are not established by the lemma itself.

\paragraph{Verification boundary.}
The accompanying standard-library script checks constructor decoding and
retraction, local-rule examples and negative controls, the exact
degree-$10$ specialization, recurrence bounds, and the exact Pell
quotient identity.  Those finite checks are reproducible corroboration.
The unrestricted mathematical statements are established by the proofs
above; no finite check establishes a consistency equivalence.


\begin{thebibliography}{99}
\raggedright
\bibitem{q41}
ProveIt, source 64, \emph{A quartic in 41 unrestricted integer
unknowns}, \src{integer41/PROOF41.md}. Pinned commit
\src{9cca523d9f27ad39d52993f4065cbd7e58f8bdd4}.
Vendored in \src{provenance/} and \src{integer_coefficients/proofs/}.
\url{https://github.com/VladimirReshetnikov/ProveIt/tree/9cca523d9f27ad39d52993f4065cbd7e58f8bdd4/Computability/DiophantineConsistency/con-zfc-diophantine-forms/64-structural-reductions}

\bibitem{q44}
ProveIt, source 64, \emph{A quartic in 44 unrestricted integer
variables}, \src{integer41/PROOF44.md}, same pinned commit.
Contains the absolute-radix sign argument and the original
source-34 coefficient binding.

\bibitem{p13old}
ProveIt, source 61, \emph{Thirteen positive unknowns, exact
degree 34, 93 operations}, \src{positive/PROOF13.md}.
\url{https://github.com/VladimirReshetnikov/ProveIt/tree/9cca523d9f27ad39d52993f4065cbd7e58f8bdd4/Computability/DiophantineConsistency/con-zfc-diophantine-forms/61-congruence-carriers}

\bibitem{p14}
ProveIt, source 61, \emph{Fourteen positive unknowns, exact
degree 34, 93 operations}, \src{positive/PROOF14.md}.
Specifies the normalized source, signed auxiliary-index
interface, and canonical compact witnesses.

\bibitem{signed}
ProveIt, inherited \src{SIGNED_DOMAIN_29.md},
\src{positive61/SIGNED_INDEX_LEMMA.md}, and
\src{PELL_RELAXED_AUXILIARY_PROOF.md}.
Exact locations and identities are recorded in the coefficient
bundles' provenance manifests.

\bibitem{compact}
ProveIt, source 44, \emph{Complement-packed native clauses:
a smaller complete coefficient tuple},
\src{native/PROOF.md}, \src{COEFFICIENT_BINDING.json}, and
\src{build_coefficients.py}.
\url{https://github.com/VladimirReshetnikov/ProveIt/tree/9cca523d9f27ad39d52993f4065cbd7e58f8bdd4/Computability/DiophantineConsistency/con-zfc-diophantine-forms/44-b4-sparse}

\bibitem{parity}
ProveIt, source 36, \src{coefficients/PROOF.md}, Section 6.
Canonical square scale, odd normalized and ordinary multipliers,
and odd cubic scale; preserved as
\src{compact_coefficients/native/source/coefficients__PROOF.md}.

\bibitem{mixed}
ProveIt, source 34, \emph{Mixed-parity unrounded coefficient
adapter}. Complete finite recipe and mathematical proof
vendored in \src{integer_coefficients/}; base exponent \(1064921\).

\bibitem{oldsextic}
ProveIt, source 62, \src{quartic42/companions},
degree-six equation in 35 unrestricted integers.
\url{https://github.com/VladimirReshetnikov/ProveIt/tree/9cca523d9f27ad39d52993f4065cbd7e58f8bdd4/Computability/DiophantineConsistency/con-zfc-diophantine-forms/62-sign-and-norm}

\bibitem{wave66}
ProveIt, source 66, \src{66-gap-congruence/README.md} and
metric receipts, same pinned commit.
\url{https://github.com/VladimirReshetnikov/ProveIt/tree/9cca523d9f27ad39d52993f4065cbd7e58f8bdd4/Computability/DiophantineConsistency/con-zfc-diophantine-forms/66-gap-congruence}

\bibitem{wave67}
ProveIt, source 67, \src{67-oriented-guards/README.md}.
\url{https://github.com/VladimirReshetnikov/ProveIt/tree/9cca523d9f27ad39d52993f4065cbd7e58f8bdd4/Computability/DiophantineConsistency/con-zfc-diophantine-forms/67-oriented-guards}

\bibitem{status}
ProveIt, \src{Computability/DiophantineConsistency/README.md}
and \src{STATUS.md}. Distinguishes closed executable Lean
theorems, native literal equations, and the internal ZFC target.
\url{https://github.com/VladimirReshetnikov/ProveIt/tree/9cca523d9f27ad39d52993f4065cbd7e58f8bdd4/Computability/DiophantineConsistency}
\end{thebibliography}
\end{document}
